\headline={\hfil\sl A Small Sample\hfil}
\footline={\hfil\folio\hfil}
\centerline{\bf Typesetting a Little Plain TeX}
\medskip
\noindent Plain \TeX\ sets paragraphs by breaking a long line of boxes and glue
into lines of equal width, choosing the breaks that make the whole paragraph
look best. It does the same for pages\footnote{$^1$}{A footnote is set at the
bottom of the page, under a short rule.} and for mathematics, where
$$\int_0^1 {dx\over 1+x^2}={\pi\over4}\qquad\hbox{and}\qquad
\sum_{k=1}^{n}k={n(n+1)\over2}.$$
\medskip
\halign{\hfil#\hfil\quad&\hfil#\hfil\quad&\hfil#\hfil\cr
\bf n&\bf $n^2$&\bf $2^n$\cr
\noalign{\smallskip}
1&1&2\cr 2&4&4\cr 3&9&8\cr 4&16&16\cr 5&25&32\cr}
\bigskip
\noindent Hyphenation is by Liang's algorithm: a paragraph of ordinary running
text such as this one, with words like typesetting, computer, hyphenation and
mathematical, lets the line breaker split them when that gives a better paragraph.
\vfill\eject
\noindent The second page holds only this line, followed by a rule.
\vskip 1in \hrule
\vfill\eject
